% 06-bai-tap.tex -- End-of-chapter exercises
%
% Scope: three tiers as per SPEC

\exercises

\index{bài tập!kiểm tra hiểu|(}%

\par\textbf{Hướng dẫn chung:} Code cho chương này nằm trong repo
\texttt{rnn-to-transformer-lab} tại tag \texttt{ch02}.
Mở file \mintinline{text}{ch02_symptoms.py}.
Mọi số liệu đến từ \mintinline{text}{python verify.py}.

\index{bài tập!kiểm tra hiểu|)}%

\tierunderstand

\begin{bt}
  \item Bài toán cộng yêu cầu mạng nhớ \textbf{hai} giá trị trong toàn bộ
    chuỗi. Nếu ta đổi thành ba vị trí đánh dấu, độ khó tăng lên thế nào?
    Giải thích bằng khái niệm dung lượng trạng thái ẩn.

  \item Nhìn vào bảng kết quả copy task. Với $T_{\text{mem}} = 20$, loss
    sau huấn luyện (8.37) cao hơn loss trước huấn luyện (7.91). Về mặt
    con số, điều này có nghĩa là gì? Mạng đang làm gì sai mà loss tăng
    lên thay vì giảm xuống?

  \item Tại sao gradient ở bước $t=1$ luôn nhỏ hơn gradient ở bước $t=T$
    trong copy task? Dùng công thức BPTT từ chương~01 để giải thích, không
    cần tính toán.

  \item Bài báo Bengio 1994 chứng minh rằng xác suất gradient được bảo
    toàn giảm theo hàm mũ của $T$. Nếu xác suất đó là $0.9^T$, thì với
    $T=50$, xác suất là bao nhiêu? So sánh với tỉ lệ norm gradient bạn
    thấy trong thí nghiệm.
\end{bt}

\tierapply

\index{bài tập!để tay vào|(}%

\begin{bt}
  \item Mở \mintinline{text}{ch02_symptoms.py}, tìm hàm
    \mintinline{python}{train_adding}. Tăng $T$ lên 100 và chạy lại với
    20000 mẫu. Loss sau huấn luyện là bao nhiêu? Tỉ lệ norm gradient thay
    đổi thế nào so với $T=50$?

  \item Sửa hàm \mintinline{python}{train_copy} để in ra loss \textbf{tại
    từng epoch} (mỗi 100 epoch). Vẽ đồ thị loss theo epoch cho $T_{\text{mem}}
    = 5, 10, 20$. Ở $T_{\text{mem}}$ nào thì loss bắt đầu tăng thay vì giảm?

  \item Thay hàm kích hoạt từ $\tanh$ sang ReLU trong cả hai bài toán
    (\mintinline{python}{np.maximum(0, a)}). Chạy lại thí nghiệm. Gradient
    còn biến mất không? Nếu không, vì sao ReLU không giải quyết được vấn
    đề này trong RNN? (Gợi ý: ReLU không có upper bound.)

  \item Dùng \mintinline{python}{torch.nn.RNN} (PyTorch) cài đặt lại bài
    toán sao chép với $T_{\text{mem}} = 20$. Để PyTorch tự tính gradient.
    In ra norm của gradient theo $\mat{W}_{hh}$ tại mỗi bước thời gian.
    Kết quả có khớp với BPTT thủ công không?
\end{bt}

\index{bài tập!để tay vào|)}%

\tierextend

\index{bài tập!đẩy xa hơn|(}%

\begin{bt}
  \item Thiết kế một bài toán tổng hợp mới kiểm tra trí nhớ dài hạn theo
    cách khác với adding và copy. Yêu cầu: (a) bài toán phải đơn giản đến
    mức mô tả được trong ba câu, (b) một RNN với $d_h = 32$ phải thất bại
    ở $T > 30$, (c) một LSTM (sẽ học ở chương~4) phải giải được ở cùng $T$.
    Cài đặt và đo.

  \item Bài báo của Hochreiter 1991 viết bằng tiếng Đức. Tìm một bản dịch
    hoặc bài tổng quan tiếng Anh, đọc phần mô tả thí nghiệm của ông, và
    so sánh với thí nghiệm trong chương này. Hochreiter dùng bài toán gì?
    Kiến trúc RNN của ông khác gì so với RNN trong chương này?

  \item (Câu hỏi mở) Chương này đo $\|\partial L / \partial \mat{W}_{hh}\|$
    như một chỉ số cho sức khỏe gradient. Nhưng liệu norm của gradient theo
    \textbf{input} ($\partial L / \partial \mat{W}_{xh}$) có bị ảnh hưởng
    bởi vanishing không? Nếu không, vì sao? Nếu có, cơ chế khác gì? Gợi ý:
    truy ngược BPTT cho $\partial L / \partial \vect{x}_t$.
\end{bt}

\index{bài tập!đẩy xa hơn|)}%
